% !PRINTED 510
\weberpage{510}

Wir können $\Phi(M,u)$ als \gesperrt{ganze Funktion von} $u$ darstellen und können dann diese Funktion durch ihren größten gemeinschaftlichen Teiler mit $H_{-D}(u)$ ersetzen.

Wir können $\alpha$ und $\beta$ positiv annehmen. Für $\beta$ liegt dies in den bereits gemachten Voraussetzungen. Für $\alpha$ können wir es immer dadurch erreichen, daß wir $B$ um ein Vielfaches von $2A$ vermehren, wodurch wir zu einer äquivalenten (parallelen) Form kommen. Dadurch kann $Bx+y$ positiv gemacht werden.

Bestimmen wir also den Quotienten
\[
\frac{\eta\left(\dfrac{\gamma+\delta\omega}{\alpha+\beta\omega}\right)}{\eta(\omega)} = E\binom{\alpha,\ \beta}{\gamma,\ \delta}
\]
nach § 38, (15), so ergibt sich, da $\beta$ gerade oder ungerade ist, je nachdem $x$ gerade oder ungerade ist:

\weq{12}{\begin{aligned}
x\equiv 0 \pmod 2:&\quad M=\left(\frac{\beta}{\alpha}\right)i^{3\frac{\alpha-1}{2}}e^{\frac{\pi i}{4}\alpha(\gamma-\beta)}\sqrt{\alpha+\beta\omega}^{\,3},\\
x\equiv 1 \pmod 2:&\quad M=\left(\frac{\alpha}{\beta}\right)i^{3\frac{1-\beta}{2}}e^{\frac{\pi i}{4}\beta(\alpha+\delta)}\sqrt{-i(\alpha+\beta\omega)}^{\,3}.
\end{aligned}}

Die Quadratwurzeln
\weq{13}{\begin{aligned}
\sqrt{\alpha+\beta\omega}&=\sqrt{\dfrac{y+x\sqrt D}{2}},\\
\sqrt{-i(\alpha+\beta\omega)}&=\sqrt{-i\dfrac{y+x\sqrt D}{2}}
\end{aligned}}
haben positiven reellen Bestandteil. Diese Werte von $M$ genügen also der Gleichung (11). Kommen unter ihnen gleiche vor, so kann der Grad der Gleichung (11) durch Absonderung mehrfacher Wurzeln auf rationalem Wege erniedrigt werden.

Da die Gleichung (11) zur Zerlegung der Klassengleichung nach den Vorzeichen von $M$ angewandt werden soll, so ist es von Wichtigkeit, zu entscheiden, ob unter den zu demselben $\omega$ gehörigen Werten von $M$ in (12) solche vorkommen, die sich nur durch das Vorzeichen unterscheiden.

Wir untersuchen, wann zwei verschiedene von den Werten (12), etwa $M$, $M'$, dieselbe 8te Potenz haben. Ist
\weq{14}{M^8=M'^8,}
so muß, wenn $\varrho$ eine zwölfte Einheitswurzel und $x,y;\ x',y'$ zwei verschiedene Lösungen der Gleichung (8) sind:
\weq{15}{y'+x'\sqrt D=\varrho(y+x\sqrt D).}
